\section{Introduction}
\label{sec:introduction}

\subsection{Background}

Bitcoin was introduced by Satoshi Nakamoto as a peer-to-peer electronic cash
system: a network in which payments can be validated by participants directly,
without requiring a financial intermediary to order transactions or prevent
double spending \cite{nakamoto-bitcoin}. Its core design combines digital
signatures, a public append-only history, a proof-of-work timestamp chain, and
local validation by independently operated nodes. Transactions spend previously
created outputs and create new outputs; blocks commit ordered transaction sets;
nodes accept the valid chain with the greatest accumulated proof of work.

The consequence is a protocol whose authority is operational rather than
institutional. No server, committee, or publication grants validity to a
transaction or block. A block is useful only if economically relevant validating
nodes accept it under their consensus rules, and a transaction is final only to
the extent that it remains in the accepted proof-of-work history. This makes
Bitcoin unusually sensitive to precise definitions: serialization, hash
preimages, script evaluation, subsidy, difficulty adjustment, lock-time,
witness commitments, and deployment rules are not merely engineering details.
They are the boundary between a node that follows the same network and a node
that has forked into a different system.

\subsection{Why this specification exists}

Bitcoin does not have a single normative protocol specification in the style of
an Internet RFC. Its design is distributed across the original paper,
Bitcoin Improvement Proposals, Bitcoin Core source code, release notes,
developer documentation, test vectors, and long-running network practice
\cite{bitcoin-bips-repo,bitcoin-core,bitcoin-core-v31,bitcoin-developer-reference}.
That distribution has served Bitcoin well: it keeps proposals reviewable, lets
implementations evolve conservatively, and avoids granting authority to a
document detached from deployed validation behavior. It also leaves developers
with a practical problem. To implement, audit, test, or reason about Bitcoin,
they must assemble one protocol from many artifacts that use different levels
of precision and often mix consensus, policy, wallet behavior, and operational
guidance.

This specification is intended to be the consolidated technical reference for
that protocol. It states Bitcoin as data structures, encodings, identifiers,
state variables, validation predicates, state transitions, and peer messages.
It is not a proposal to change Bitcoin, not a replacement for BIPs, and not an
authority over consensus. Where this document conflicts with economically
relevant node validation, node validation wins. Its purpose is to make the
protocol easier to inspect: one place where the rules needed by independent
implementers, reviewers, researchers, and protocol engineers are expressed
coherently.

\subsection{Related work}

The original Bitcoin paper defines the motivating problem and the high-level
mechanism: timestamped proof of work over a chain of transaction-containing
blocks \cite{nakamoto-bitcoin}. BIPs refine that design by documenting concrete
changes, conventions, and interoperability standards, including consensus soft
forks, peer services, wallet formats, and application-level conventions
\cite{bip-0123,bitcoin-bips-repo}. Bitcoin Core is the dominant full-node
implementation and the principal executable reference for contemporary
validation and relay behavior \cite{bitcoin-core,bitcoin-core-v31}.
Developer references and generated API documentation are valuable guides, but
they are not a complete formal specification of the network protocol
\cite{bitcoin-developer-reference}.

Other cryptocurrency systems have benefited from standalone protocol
specifications. The Zcash protocol specification separates high-level protocol
concepts from concrete encodings and upgrade-specific consensus rules
\cite{zcash-protocol-spec}. Ethereum's Yellow Paper presents Ethereum as a
transaction-based state machine with explicit state-transition functions
\cite{ethereum-yellowpaper}. The JAM graypaper follows the same tradition of a
single coherent formal model for an independently implementable protocol
\cite{graypaper-repo}. This document applies that style to Bitcoin while
respecting Bitcoin's different governance reality: no paper, repository, or
committee has authority to redefine the deployed protocol.

\subsection{Methodology}

The baseline for this draft is Bitcoin Core source, its
implemented-BIPs index, and the BIPs cited throughout this document are treated
as evidence for contemporary full-node behavior
\cite{bitcoin-core-v31}. The specification
separates three classes of behavior:

\begin{enumerate}[leftmargin=*]
  \item Consensus rules: predicates that determine whether blocks and
        transactions can be accepted into the active chain.
  \item Peer and relay rules: network records and local admission behavior used
        by nodes to exchange blocks, transactions, addresses, and negotiation
        state.
  \item Local implementation behavior: storage layout, indexes, caches, RPC
        surfaces, wallet policy, user interface behavior, and performance
        choices that may explain Bitcoin Core but are not themselves protocol
        requirements.
\end{enumerate}

Only the first two classes belong in the main specification. Implementation
details are included only when they expose a protocol boundary or prevent an
ambiguous reading of consensus or peer behavior. This keeps the document useful
as a source-of-truth reference without turning it into a commentary on one code
base.
